% Supplementary material, placed after the references as the journal's instructions ask. Sections S1 to S8 are referred
% to from the main text by their numbers.
\renewcommand{\thesection}{S\arabic{section}}
\setcounter{section}{0}

\section{Conditions, status of the analyses and changes}
\label{app:status}

\paragraph{Order of the computations.} Every computation ran under a written condition committed to the author's version
control before it ran. The conditions are in the archive, in Japanese, with an English summary. They are records in
the author's repository, not third-party registrations: their order is shown by the repository, not by an independent
time stamp.

\begin{center}\footnotesize
\begin{tabularx}{\textwidth}{@{}c l X@{}}
\toprule
Order & Status & Computation and what was fixed before it \\
\midrule
1 & [P] & Primary: the 90 definitions, rates, ranks, $\varphi$, agreement in time, the decomposition of counts and lengths
for 9 definitions, an estimator-level comparison (in the archive); the criteria for ``moves'' (25\% of definitions) and for
overdispersion \\
2 & [A] & Length decomposition for all 90 definitions; the criterion that the stall-to-stall share stays the largest \\
3 & [A] & Intervals for rates, ratios and $\varphi$; the KLiCKe score association; the criterion for a reversal of
the score association (one definition with an interval on the other side of zero) and, set after the reversed point
values had been seen and before the intervals, for a ratio (25\% of definitions for a pair) \\
4 & [E] & Partial association given the key-press rate (withdrawn as a covariate, below) \\
5 & [E] & Partial association given typing speed (median interval $< 2$~s) \\
6 & [E] & Six ways of measuring speed and the within-prompt adjustment; the reading sentence stated in
Section~\ref{sec:res-adjust} \\
7 & [A] & After a first internal pre-submission check: the pause family (36 definitions), intervals with the writer
quantile estimated again, and the adjustment for essay and observation length; the readings of each, including that
the two-stage intervals replace the earlier ones regardless of which is narrower \\
8 & [A] & The restriction of all 90 definitions to the time after each unit's 100th interval (S2); its readings \\
9 & [A] & After a second internal check: two-stage intervals with the writer quantile estimated causally (S2), the
counting rule that a reversal needs both the point value and the interval on the other side, and the reading that the
new intervals replace the one-stage ones only if at most 5\% of point values fall outside their intervals; 27\% of
the $\varphi$ values did, so the primary intervals returned to the one-stage ones (the score association,
positive for 10 definitions under the one-stage intervals, had been positive for 14 under those of order 7) \\
10 & [D] & After a third internal check, no new computation: the reading by classes of interval, with imprecise intervals
as a class of their own, and the grid of default and precedented levels, counted from the existing intervals (in the
archive) \\
11 & [A] & After a fourth internal check: the share of exposure in the first $W$ seconds of each observation interval,
and all 90 definitions with each interval cut to start 60~s after its start and not before the 100th interval (S2);
its readings \\
\bottomrule
\end{tabularx}
\end{center}

\paragraph{Status of each conclusion.} [P] primary, quantity and criterion committed before any result; [A] added
after the primary results, quantity and criterion committed before that quantity was seen; [E] exploratory, chosen
after seeing results; [D] descriptive, counted after seeing results.

\begin{center}\footnotesize
\begin{tabularx}{\textwidth}{@{}X l@{}}
\toprule
Conclusion & Status \\
\midrule
Direction of the stall--score and pause--score association changes (criterion: one definition on the other side of
zero). For pauses the reading was stated over all 36 definitions and was met there, with the control supplying 18
positive intervals; the main text reports the 18 fixed-threshold definitions, in which it is also met (3 positive and
7 negative intervals), and treats the 18 writer-quantile ones as a control. Presence of the pause association: by the stall criterion applied after seeing the results, it moves on the
18 (56\% of intervals exclude zero) and not on the 36 (78\%) & [A] \\
Reading by classes of interval (90\% of classified definitions in one class; imprecise intervals counted apart) and
the grid of default and precedented levels & [D] \\
Presence of the association does not move (criterion: 25--75\% of intervals excluding zero) & [A] \\
Strength of the association (negligible to small) & [D] \\
Given typing speed, negative for all stall definitions & [E] \\
Given essay length, negative for none & [A] \\
Rank order of datasets changes (criterion: 25\% of definitions) & [P] \\
Rate ratios reverse, point value and interval on the other side of 1 (criterion set after seeing reversed point values) & [D] \\
Within-writer $\varphi$ reverses at 1.5 (in 2 of the 4 CS1 cells over 25\%) & [P, A] \\
CS1 stall-to-stall share stays the largest & [A] \\
\bottomrule
\end{tabularx}
\end{center}

\paragraph{Why the covariate was chosen twice.} After the score association was computed, we asked whether the stall
rate of writer-quantile definitions merely counted typing. The first covariate, the key-press rate, divides presses by
a time that contains the stalls themselves, so adjusting for it removes part of the stall rate; we withdrew it as
evidence and report it only for reference (S3). The second, typing speed measured on short intervals, does not contain
the stall labels. Because the result of the second changed the direction of the association, we varied the measure of
speed six ways under a condition fixed before that computation, and treat all adjusted results as exploratory.

\paragraph{Departures from the conditions.} (1) The exposure of the count model entered as a fixed slope, because the
variational mixed model has no offset; the slopes were 0.99 (CS1 2019), 1.04 (CS1 2021) and 1.29 (CoAuthor). (2) The
count model used the L-BFGS-B optimiser; the default did not finish on the real data, and the two agreed to two
decimals on simulated data. (3) The exposure of a CS1 work block is the overlap of the observation intervals with the
time from the start of the block to the start of the next, so that stalls between blocks do not fall into zero
exposure. (4) The recovery check of the decompositions ran in the same execution as the estimation; the judgement used
only recovered levels, so the order does not change it. (5) The writer-quantile 0.995 variant of the reference
implementation had been built with a minimum of 200 intervals; with that minimum the onsets agreed in all units. (6) The
resumption of a stall is restricted to the same observation interval; a first version took the next character across
intervals and was rerun. (7) Writers with zero stalls are left out of the writer-expectation $\varphi$. (8) The
10-minute minimum of observation was applied to all datasets. (9) The prior SD of the variance components of the count
model was set to 2. (10) The overdispersion criterion was counted over the five cells. (11) Agreement is counted as 0
when one of the two definitions has no stall. (12) After a two-system code review, the selftests and 20 mutants were
completed (all mutants fail the checks, the unmodified code passes). (13) The length model took the maximum REML
likelihood over three optimisers (S5). (14) The variational count model gave different values on repeated runs (S6).
(15) The text column of the KLiCKe score file was read with replacement of undecodable characters; the identifier,
prompt and score columns are ASCII, and essay length counts characters of that column. (16) The result files were
rewritten after a conversion of integer types, with the same seed. (17) The check that the one-stage intervals of the
two-stage code reproduce the earlier intervals within 0.05 failed at 0.070; the largest differences were upper ends of
quantities involving CoAuthor (41 writers), and two runs of the same code with different seeds differed by 0.056, so we
widened the tolerance to 0.08 (3 of 864 comparisons exceeded 0.05). (18) A summary count of the earlier intervals by
threshold family was corrected before the results were read. (19) A synthetic control of the two-stage code was
rebuilt before the run because its constant intervals included the long gap. (20) The first mutation check ran in a
directory where imports failed; it was repeated in a copy of the layout, and controls were added for the mutants that
survived.

\paragraph{Holdouts.} The CS1 2019 holdout of 146 writers was drawn at random before any analysis of these data and
opened after the companion study's confirmatory test (145 had key presses). The KLiCKe holdout of 1,498 writers had
been opened in the same way. The CS1 2021 holdout (13 writers) and the CoAuthor holdout (18 writers, 288 sessions)
were not opened; checks of the work files confirmed that none of their units entered the analyses. The rule of 10
minutes of observation (the time within the observation intervals, not the time on the KLiCKe timer) excluded 47 CS1
2019 assignments, 84 KLiCKe essays, 2 CoAuthor sessions and none in CS1 2021.

\section{Detector, conversion and intervals}
\label{app:detector}

\paragraph{Conversion of CS1.} An insertion of one character with a recorded key became a character press, a deletion
with a recorded key a deletion press. Insertions of several characters (mostly pastes) and edits without a recorded key
did not become key presses. The mapping is that of the companion study \cite{miura2026runup}.

\paragraph{Threshold and evaluation times.} The writer quantile is the nearest-rank quantile of the inter-keystroke
intervals within the observation intervals of the unit so far, updated causally, defined once 100 intervals exist.
Fixed thresholds do not depend on the writer. Condition A is evaluated at the union of the key-press times and the
key-press times plus $W$, each time once; the resumption is the first character press after $t_0$, so a length is
always positive.

\paragraph{Observation intervals.} The default intervals are where the companion study's estimator, re-run on the key
stream with a release added 100~ms after each press, wrote per-second rows, split where rows are more than 2~s apart.
They cover a median of 99.9\% of the time from the first to the last key press of a KLiCKe essay (1st percentile
80\%) and of a CoAuthor session from its start; holes are mostly 3~s long and contain a key press after a gap of
several seconds (median 7~s in KLiCKe). A CS1 unit is replayed work block by work block (blocks split at 30-minute
gaps), so time between blocks lies outside. The two other segmentations use key times alone.

\paragraph{Time before the 100th interval.} The share of exposure before each unit's 100th interval in the default
segmentation was 3.4\% (CS1 2019), 5.0\% (CS1 2021), 3.2\% (KLiCKe) and 10.3\% (CoAuthor), and 91, 14, 0 and 1 units
never reached 100 intervals (2.8--8.4\% with key-press gaps of 5~min). Restricting all 90 definitions to the time after the
100th interval of each unit (both numerator and exposure; one-stage intervals, compared with the one-stage values of
the full time) changed the baseline stall--score correlation from $-0.063$ to $-0.081$ [$-0.108$, $-0.053$]; 9
definitions instead of 10 had a positive interval, 74 a negative one, and no interval changed from one side of zero to
the other. The rank order differed for 80 definitions, as before. The rate ratios reversed for 22 definitions between
the two CS1 editions (20 before; 19 with a fixed $\theta$, 18 before), 66, 53, 8 and 2 for the other pairs (66, 53, 6
and 0 before). The reversals of $\varphi$ at 1.5, counted by the rule of
the main text, were 17 (CS1 2019 blocks), 57 (CS1 2019 assignments), not counted (CS1 2021 blocks), 26 (CS1 2021
assignments) and 0 (CoAuthor); in the full time they were not counted, 54, 49, 19 and 0. In the restricted time the
baseline interval of the CS1 2019 blocks, [1.28, 1.45], no longer contains 1.5, whereas that of the CS1 2021 blocks,
[1.37, 1.94], does. The CS1 2021 assignments newly crossed 25\% of the definitions, and the CS1 2019 assignments stayed above it.

\paragraph{The first seconds of each observation interval.} Condition A is evaluated only from $W$ after the start of
an observation interval, so no definition can detect a stall in the first $W$ seconds of an interval. These seconds
held 1.1--2.8\% (15~s), 2.1--5.6\% (30~s) and 4.0--10.9\% (60~s) of the exposure in the default segmentation, and
0.5--9.2\% in the two others, by dataset and window. We cut each interval $[a, b]$ to $[\max(a + 60\,\mathrm{s}, w), b]$,
with $w$ the time of the unit's 100th interval, and dropped empty intervals, for all 90 definitions (exposure on the
cut intervals, stalls counted when $t_0$ lay in them, the 5-minute rule applied to the cut exposure, one-stage
intervals). The baseline stall--score correlation was $-0.073$ [$-0.100$, $-0.045$]; 9 definitions had a positive
interval (all with a 15-s window and a writer quantile; 10 in the full time), 74 a negative one (72), and no interval
changed from one side of zero to the other. The rank order differed for 81 definitions (80). The rate ratios reversed
for 23 definitions between the two CS1 editions (20; 19 with a fixed $\theta$, 18 before), so this pair crossed 25\%
of the definitions, and for 66, 53, 8 and 2 in the other pairs (66, 53, 6 and 0). The reversals of $\varphi$ at 1.5
were 17, 57, not counted, 24 and 0 in the order of the previous paragraph (not counted, 54, 49, 19 and 0), the same
pattern as after the 100th interval. The primary values remain those of the full time.

\paragraph{Checks.} The onsets $(s, t_0)$ of the baseline agreed with those of the reference implementation in all
7,709 units. For the window variants of 15 and 60~s and the 0.995 quantile variant (with the minimum of 200 intervals of
its reference build), the onsets agreed in all 1,078 units of the seen CS1 2019 writers. Positive controls on synthetic
streams (a stream with one stall, none, and one raised by repeated deletions) gave the expected counts for each
setting, and 20 mutants of the condition, the quantile, the segmentation and the analysis code all failed the checks.

\paragraph{Intervals.} Writers were resampled with replacement 2,000 times within each dataset with a fixed seed; the
rate of a resample is the sum of counts over the sum of exposures of its units (units with at least 5~min of
exposure). These one-stage intervals are the primary ones. For the 54 writer-quantile stall definitions (the baseline included) we also
computed two two-stage intervals. In the first, the key-press times of each unit were kept and the values of its
intervals were drawn again with replacement from all of the unit's intervals; $\theta$ was recomputed by the causal
rule and the stalls were counted again. 200 such draws were made per unit, and each writer drawn in an outer resample
took one at random. Without the inner draw, the counts agreed exactly with the detector in every block of every unit.
These intervals were about as wide as the one-stage ones (median ratio of widths for the score association 1.00, range
0.95--1.06) but shifted against the point values: 18 of 54 score values, 11 of 324 ratios and 47 of 270 values of
$\varphi$ fell outside their intervals, and under the baseline the CS1 2019 to KLiCKe ratio interval [0.807, 0.861] no
longer contained the point value 0.865. We read the shift as the effect of drawing early intervals from the whole unit,
where the causal rule uses only the intervals so far.

In the second, set after the second internal check, each interval of a unit received a Poisson(1) weight, and $\theta$ at each
time was the weighted quantile of the intervals before it, so that the draw is causal as well; again 200 draws per unit.
With unit weights it reproduced the detector's counts in every block of every unit. Its intervals had the same width
as the one-stage ones (median ratio 1.00), and all 54 score values and 324 ratios lay inside them, but 74 of 270 values
of $\varphi$ lay below their intervals (CS1 2019 blocks 33, CS1 2019 assignments 32, CS1 2021 blocks 5, CoAuthor 3,
CS1 2021 assignments 1): the redrawn $\theta$ adds its own variation to the counts within a writer and so raises
$\varphi$. By the reading set before the computation (more than 5\% outside in any quantity), the primary intervals
returned to the one-stage ones. Under the second two-stage intervals, counting a reversal only when the point value and the interval
are on the other side, the score association was positive for the same 10 definitions and negative for 74; the ratio
reversals were 20 (the two CS1 editions), 66 (CS1 2019 against KLiCKe), 53 (CS1 2021 against KLiCKe) and 7 (CS1 2019
against CoAuthor); and the reversals of $\varphi$ at 1.5 were 42 (CS1 2019 assignments), 39 (CS1 2021 blocks) and 18
(CS1 2021 assignments). Under the first two-stage intervals, 14 definitions had a positive score interval.

\section{All definitions}
\label{app:all}

Table~\ref{tab:s3} lists, for each definition, the pooled stall rates, the rank of the datasets, two rate ratios, the
score association with and without adjustment for speed (median interval $< 2$~s) and for essay and observation
length, and the share of KLiCKe writers with no stall. Table~\ref{tab:spause} gives the pause family. For stalls, the rate ratio of CS1 2019 against KLiCKe (baseline 0.86) reversed for 66
definitions (36 with a fixed $\theta$, 30 with a writer quantile), CS1 2021 against KLiCKe for 53 (36 and 17), CS1 2019
against CoAuthor for 6 (all writer quantile) and the two writing datasets for none; the baseline interval of CS1 2021
against CoAuthor included 1, so its reversals were not counted. In the pause family the rate ratio reversed for no definition between
the two CS1 editions, for 32 between CS1 2019 and CoAuthor (14 fixed, 18 quantile), for 31 between CoAuthor and KLiCKe
(13 and 18), and for 5 and 3 between the CS1 editions and KLiCKe (all fixed); by the criterion set after seeing point
values, the direction of a ratio changed with the definition in the pause family too. The rank order
of the four datasets differed from the baseline order for 80 of 89 definitions; the baseline rates of CS1 2019 and
KLiCKe are close, and their swap alone accounts for 11 of the 80. Within the four prompts of KLiCKe, as point values only, the median of the
prompt correlations between stall rate and score was negative for 74 definitions and positive for 15. Among the 10
definitions with a positive score interval, the share of KLiCKe writers with no stall was 0.1--1.3\%, against
0.1--35\% across all 89, so the positive corner is not a floor effect of rare stalls. The definitions that lowered
$\varphi$ most split the observation at key-press gaps of 5~min or used a fixed 30-s threshold. For pauses,
the number of definitions whose $\varphi$ interval lay on the other side of 1.5 from the baseline stall $\varphi$ was 26
(CS1 2019 blocks), 1 (CS1 2019 assignments), 8 (CS1 2021 blocks), 1 (CS1 2021 assignments) and 6 (CoAuthor sessions);
we report these counts without a reading.

\paragraph{Ways of measuring speed [E].} The partial correlations of score and stall rate given speed, as (negative /
including zero / positive) out of 89, with their range: maximum interval 1~s, median, 89/0/0, $-0.28$ to $-0.03$;
2~s, median, 89/0/0, $-0.29$ to $-0.03$; 3~s, median, 89/0/0, $-0.29$ to $-0.04$; 2~s, mean, 89/0/0, $-0.30$ to $-0.05$;
2~s, intervals between two character presses only, 87/2/0, $-0.29$ to $-0.02$; median of 10--750~ms (the measure of
\citeN{koutcheme2026keystroke}), 89/0/0, $-0.28$ to $-0.03$; 2~s, median, within prompts (weighted mean over the four
prompts), 89/0/0, $-0.28$ to $-0.03$. Under the baseline the partial correlation was $-0.17$ to $-0.19$ for every
measure. The correlation of the stall rate with speed ranged from $-0.35$ to $+0.07$ over the 89 definitions. The
reading set in advance, that the direction would ``also depend on how speed is measured'' if any interval included zero
or was positive, was met by the 2 intervals that included zero; by the counting rule of the main text no direction
reversed.

\paragraph{Key-press rate (reference only).} The score correlated with the key-press rate at 0.62 [0.60, 0.63]. Given
the key-press rate, the partial correlation of score and stall rate ranged from $-0.054$ to $+0.106$, positive with an
interval excluding zero for 48 definitions and negative for 27; the sign split by condition A. Because the key-press
rate shares its denominator with the stalls, we do not use these values.

\paragraph{Agreement in time.} Two labels agree on a stall when the onset of one falls between the start and the
resumption of a stall of the other. For each definition and dataset we computed the share of baseline onsets that fall
within a stall of the definition, and the reverse. The smaller of the two was below 0.5 for 51\% of the definitions in CS1 2019, 36\%
in CS1 2021, 77\% in KLiCKe and 80\% in CoAuthor. Among the 63 pairs of a definition and a dataset whose rate was within
20\% of the baseline rate, 11 had the smaller agreement below 0.5 (4 in CS1 2019, 6 in KLiCKe, 1 in CS1 2021); for
example, in KLiCKe a definition with $W = 15$~s, the characters-only rule and a fixed 10-s threshold had a rate within
6\% of the baseline and an agreement of 0.16--0.17. The agreement is not corrected for chance.

{\scriptsize
\setlength{\tabcolsep}{2.5pt}
\begin{longtable}{@{}l r l l r r r r l l l l r r r@{}}
\caption{All definitions. $\theta$: writer quantile or fixed seconds; A: net or characters only; Obs.: default
intervals or key-press gaps of 5 or 30~min. Rates per hour for CS1 2019, CS1 2021, CoAuthor and KLiCKe. Rank: highest
rate first (a CS1 2019, b CS1 2021, c CoAuthor, k KLiCKe). Ratios with 95\% intervals. $\rho$: score association with
95\% interval; adj.: partial correlation given speed [E]; len.: given essay and
observation length. Zero: \% of KLiCKe writers with no stall. Baseline marked
with an asterisk.}\label{tab:s3}\\
\toprule
$\theta$ & $W$ & A & Obs. & CS19 & CS21 & CoA & KL & Rank & CS19/KL & CS19/CS21 & $\rho$ & adj. & len. & Zero \\
\midrule
\endfirsthead
\toprule
$\theta$ & $W$ & A & Obs. & CS19 & CS21 & CoA & KL & Rank & CS19/KL & CS19/CS21 & $\rho$ & adj. & len. & Zero \\
\midrule
\endhead
\bottomrule
\endfoot
% Generated from the aggregate result files by the authors' script; do not edit by hand.
0.95 & 15 & net & def. & 36.4 & 26.8 & 36.6 & 40.6 & kcab & 0.90 [0.87, 0.92] & 1.36 [1.26, 1.47] & 0.026 [$-$0.00, 0.05] & $-$0.061 & 0.125 & 0.1 \\
0.95 & 15 & net & 5 & 46.0 & 39.7 & 37.3 & 42.3 & akbc & 1.09 [1.07, 1.11] & 1.16 [1.11, 1.21] & 0.010 [$-$0.02, 0.04] & $-$0.096 & 0.135 & 0.2 \\
0.95 & 15 & net & 30 & 36.1 & 26.1 & 37.3 & 40.6 & kcab & 0.89 [0.86, 0.91] & 1.38 [1.28, 1.50] & 0.028 [0.00, 0.05] & $-$0.056 & 0.123 & 0.1 \\
0.95 & 15 & char & def. & 27.2 & 21.9 & 27.1 & 25.0 & ackb & 1.09 [1.07, 1.11] & 1.24 [1.17, 1.33] & $-$0.070 [$-$0.10, $-$0.04] & $-$0.087 & 0.111 & 0.4 \\
0.95 & 15 & char & 5 & 34.0 & 32.0 & 27.7 & 25.9 & abck & 1.31 [1.29, 1.33] & 1.06 [1.03, 1.10] & $-$0.087 [$-$0.12, $-$0.06] & $-$0.119 & 0.117 & 0.5 \\
0.95 & 15 & char & 30 & 27.0 & 21.3 & 27.7 & 25.0 & cakb & 1.08 [1.06, 1.10] & 1.27 [1.18, 1.35] & $-$0.067 [$-$0.10, $-$0.04] & $-$0.083 & 0.109 & 0.3 \\
0.95 & 30 & net & def. & 25.4 & 20.6 & 11.7 & 21.7 & akbc & 1.17 [1.14, 1.21] & 1.24 [1.15, 1.32] & $-$0.104 [$-$0.13, $-$0.08] & $-$0.182 & 0.086 & 0.9 \\
0.95 & 30 & net & 5 & 31.6 & 29.8 & 12.0 & 22.5 & abkc & 1.40 [1.37, 1.44] & 1.06 [1.02, 1.11] & $-$0.118 [$-$0.14, $-$0.09] & $-$0.204 & 0.090 & 1.4 \\
0.95 & 30 & net & 30 & 25.3 & 20.1 & 11.9 & 21.8 & akbc & 1.16 [1.13, 1.19] & 1.25 [1.17, 1.34] & $-$0.104 [$-$0.13, $-$0.08] & $-$0.178 & 0.085 & 0.8 \\
0.95 & 30 & char & def. & 16.8 & 15.1 & 7.1 & 11.4 & abkc & 1.47 [1.44, 1.50] & 1.11 [1.06, 1.17] & $-$0.159 [$-$0.19, $-$0.13] & $-$0.184 & 0.048 & 3.0 \\
0.95 & 30 & char & 5 & 20.4 & 21.4 & 7.2 & 11.8 & bakc & 1.74 [1.70, 1.77] & 0.95 [0.93, 0.98] & $-$0.168 [$-$0.20, $-$0.14] & $-$0.200 & 0.048 & 3.6 \\
0.95 & 30 & char & 30 & 16.7 & 14.8 & 7.2 & 11.5 & abkc & 1.45 [1.42, 1.48] & 1.12 [1.07, 1.19] & $-$0.157 [$-$0.18, $-$0.13] & $-$0.182 & 0.047 & 2.8 \\
0.95 & 60 & net & def. & 16.2 & 14.4 & 2.0 & 11.0 & abkc & 1.48 [1.43, 1.54] & 1.13 [1.06, 1.20] & $-$0.171 [$-$0.20, $-$0.14] & $-$0.243 & 0.062 & 7.4 \\
0.95 & 60 & net & 5 & 19.6 & 20.2 & 2.1 & 11.3 & bakc & 1.74 [1.68, 1.81] & 0.97 [0.93, 1.02] & $-$0.174 [$-$0.20, $-$0.15] & $-$0.244 & 0.061 & 9.0 \\
0.95 & 60 & net & 30 & 16.3 & 14.3 & 2.1 & 11.1 & abkc & 1.47 [1.41, 1.52] & 1.14 [1.07, 1.21] & $-$0.177 [$-$0.20, $-$0.15] & $-$0.247 & 0.059 & 6.8 \\
0.95 & 60 & char & def. & 9.0 & 9.1 & 0.9 & 4.8 & bakc & 1.88 [1.83, 1.93] & 0.98 [0.95, 1.02] & $-$0.168 [$-$0.19, $-$0.14] & $-$0.220 & 0.022 & 18.2 \\
0.95 & 60 & char & 5 & 10.2 & 12.0 & 0.9 & 4.8 & bakc & 2.14 [2.08, 2.21] & 0.86 [0.82, 0.89] & $-$0.164 [$-$0.19, $-$0.14] & $-$0.214 & 0.021 & 20.7 \\
0.95 & 60 & char & 30 & 9.0 & 9.0 & 1.0 & 4.9 & bakc & 1.85 [1.79, 1.90] & 0.99 [0.96, 1.04] & $-$0.170 [$-$0.20, $-$0.14] & $-$0.224 & 0.020 & 17.3 \\
0.99 & 15 & net & def. & 16.5 & 10.0 & 21.7 & 26.3 & kcab & 0.63 [0.60, 0.65] & 1.65 [1.50, 1.83] & 0.092 [0.06, 0.12] & $-$0.058 & 0.085 & 0.2 \\
0.99 & 15 & net & 5 & 20.7 & 14.7 & 21.8 & 27.2 & kcab & 0.76 [0.74, 0.78] & 1.40 [1.29, 1.52] & 0.077 [0.05, 0.11] & $-$0.095 & 0.092 & 0.3 \\
0.99 & 15 & net & 30 & 16.1 & 9.5 & 21.8 & 26.3 & kcab & 0.61 [0.59, 0.64] & 1.70 [1.53, 1.89] & 0.096 [0.07, 0.12] & $-$0.051 & 0.085 & 0.2 \\
0.99 & 15 & char & def. & 14.1 & 9.3 & 17.2 & 18.2 & kcab & 0.77 [0.75, 0.80] & 1.52 [1.39, 1.67] & 0.012 [$-$0.02, 0.04] & $-$0.084 & 0.066 & 0.7 \\
0.99 & 15 & char & 5 & 17.4 & 13.4 & 17.3 & 18.8 & kacb & 0.93 [0.91, 0.95] & 1.30 [1.21, 1.40] & $-$0.003 [$-$0.03, 0.02] & $-$0.115 & 0.069 & 0.8 \\
0.99 & 15 & char & 30 & 13.7 & 8.8 & 17.3 & 18.2 & kcab & 0.76 [0.73, 0.78] & 1.56 [1.43, 1.72] & 0.016 [$-$0.01, 0.04] & $-$0.079 & 0.065 & 0.6 \\
0.99$^\ast$ & 30 & net & def. & 13.1 & 8.8 & 7.1 & 15.1 & kabc & 0.86 [0.84, 0.89] & 1.49 [1.36, 1.64] & $-$0.063 [$-$0.09, $-$0.03] & $-$0.179 & 0.057 & 1.5 \\
0.99 & 30 & net & 5 & 16.1 & 12.6 & 7.2 & 15.6 & akbc & 1.03 [1.01, 1.06] & 1.28 [1.19, 1.39] & $-$0.074 [$-$0.10, $-$0.05] & $-$0.197 & 0.059 & 2.1 \\
0.99 & 30 & net & 30 & 12.9 & 8.4 & 7.2 & 15.2 & kabc & 0.85 [0.82, 0.88] & 1.53 [1.40, 1.71] & $-$0.061 [$-$0.09, $-$0.03] & $-$0.174 & 0.056 & 1.4 \\
0.99 & 30 & char & def. & 10.0 & 7.5 & 4.8 & 9.4 & akbc & 1.07 [1.04, 1.10] & 1.34 [1.24, 1.44] & $-$0.096 [$-$0.12, $-$0.07] & $-$0.169 & 0.031 & 4.6 \\
0.99 & 30 & char & 5 & 12.0 & 10.2 & 4.9 & 9.6 & abkc & 1.25 [1.22, 1.29] & 1.17 [1.10, 1.25] & $-$0.102 [$-$0.13, $-$0.07] & $-$0.179 & 0.030 & 5.5 \\
0.99 & 30 & char & 30 & 9.9 & 7.2 & 4.9 & 9.4 & akbc & 1.05 [1.02, 1.08] & 1.38 [1.28, 1.49] & $-$0.093 [$-$0.12, $-$0.06] & $-$0.165 & 0.029 & 4.4 \\
0.99 & 60 & net & def. & 9.3 & 7.2 & 1.3 & 8.0 & akbc & 1.16 [1.12, 1.21] & 1.30 [1.19, 1.43] & $-$0.138 [$-$0.17, $-$0.11] & $-$0.228 & 0.047 & 10.5 \\
0.99 & 60 & net & 5 & 11.0 & 9.8 & 1.3 & 8.2 & abkc & 1.35 [1.30, 1.40] & 1.12 [1.03, 1.23] & $-$0.139 [$-$0.17, $-$0.11] & $-$0.226 & 0.046 & 12.4 \\
0.99 & 60 & net & 30 & 9.3 & 7.0 & 1.3 & 8.1 & akbc & 1.14 [1.10, 1.19] & 1.32 [1.21, 1.46] & $-$0.143 [$-$0.17, $-$0.12] & $-$0.232 & 0.043 & 9.8 \\
0.99 & 60 & char & def. & 6.3 & 5.4 & 0.6 & 4.3 & abkc & 1.47 [1.42, 1.52] & 1.15 [1.09, 1.23] & $-$0.121 [$-$0.15, $-$0.09] & $-$0.198 & 0.022 & 21.5 \\
0.99 & 60 & char & 5 & 6.8 & 6.6 & 0.7 & 4.2 & abkc & 1.62 [1.56, 1.67] & 1.03 [0.97, 1.10] & $-$0.114 [$-$0.14, $-$0.09] & $-$0.188 & 0.021 & 24.1 \\
0.99 & 60 & char & 30 & 6.2 & 5.2 & 0.7 & 4.3 & abkc & 1.43 [1.38, 1.47] & 1.18 [1.11, 1.26] & $-$0.121 [$-$0.15, $-$0.09] & $-$0.200 & 0.021 & 20.5 \\
0.995 & 15 & net & def. & 9.7 & 5.7 & 13.9 & 19.9 & kcab & 0.49 [0.47, 0.50] & 1.71 [1.53, 1.92] & 0.137 [0.11, 0.16] & $-$0.041 & 0.058 & 0.4 \\
0.995 & 15 & net & 5 & 12.2 & 8.6 & 14.0 & 20.6 & kcab & 0.59 [0.57, 0.61] & 1.42 [1.29, 1.55] & 0.123 [0.10, 0.15] & $-$0.077 & 0.061 & 0.5 \\
0.995 & 15 & net & 30 & 9.4 & 5.4 & 14.0 & 19.9 & kcab & 0.47 [0.46, 0.49] & 1.75 [1.55, 1.97] & 0.144 [0.12, 0.17] & $-$0.034 & 0.056 & 0.4 \\
0.995 & 15 & char & def. & 8.7 & 5.4 & 11.2 & 14.3 & kcab & 0.61 [0.59, 0.63] & 1.62 [1.45, 1.81] & 0.066 [0.04, 0.09] & $-$0.067 & 0.044 & 1.1 \\
0.995 & 15 & char & 5 & 10.8 & 8.0 & 11.2 & 14.7 & kcab & 0.73 [0.71, 0.75] & 1.35 [1.24, 1.47] & 0.056 [0.03, 0.08] & $-$0.094 & 0.046 & 1.3 \\
0.995 & 15 & char & 30 & 8.5 & 5.1 & 11.1 & 14.3 & kcab & 0.59 [0.57, 0.61] & 1.67 [1.51, 1.88] & 0.075 [0.05, 0.10] & $-$0.060 & 0.043 & 1.0 \\
0.995 & 30 & net & def. & 8.1 & 5.1 & 4.5 & 12.0 & kabc & 0.67 [0.65, 0.70] & 1.58 [1.42, 1.77] & $-$0.026 [$-$0.05, 0.00] & $-$0.157 & 0.043 & 2.6 \\
0.995 & 30 & net & 5 & 9.9 & 7.6 & 4.6 & 12.3 & kabc & 0.81 [0.78, 0.83] & 1.32 [1.20, 1.45] & $-$0.034 [$-$0.06, $-$0.01] & $-$0.172 & 0.047 & 3.2 \\
0.995 & 30 & net & 30 & 7.9 & 4.9 & 4.6 & 12.1 & kabc & 0.66 [0.63, 0.68] & 1.61 [1.45, 1.80] & $-$0.023 [$-$0.05, 0.01] & $-$0.153 & 0.043 & 2.3 \\
0.995 & 30 & char & def. & 6.6 & 4.6 & 3.1 & 7.8 & kabc & 0.85 [0.83, 0.88] & 1.45 [1.32, 1.59] & $-$0.053 [$-$0.08, $-$0.02] & $-$0.151 & 0.020 & 6.7 \\
0.995 & 30 & char & 5 & 7.9 & 6.4 & 3.1 & 7.9 & kabc & 1.00 [0.97, 1.02] & 1.23 [1.14, 1.34] & $-$0.057 [$-$0.08, $-$0.03] & $-$0.157 & 0.022 & 7.8 \\
0.995 & 30 & char & 30 & 6.5 & 4.4 & 3.1 & 7.8 & kabc & 0.83 [0.81, 0.86] & 1.48 [1.35, 1.63] & $-$0.050 [$-$0.08, $-$0.02] & $-$0.147 & 0.020 & 6.3 \\
0.995 & 60 & net & def. & 6.3 & 4.4 & 0.9 & 6.6 & kabc & 0.95 [0.91, 0.99] & 1.43 [1.29, 1.58] & $-$0.112 [$-$0.14, $-$0.08] & $-$0.217 & 0.040 & 13.1 \\
0.995 & 60 & net & 5 & 7.3 & 6.1 & 1.0 & 6.7 & akbc & 1.10 [1.06, 1.14] & 1.21 [1.10, 1.34] & $-$0.109 [$-$0.14, $-$0.08] & $-$0.209 & 0.040 & 15.2 \\
0.995 & 60 & net & 30 & 6.2 & 4.3 & 1.0 & 6.7 & kabc & 0.92 [0.89, 0.96] & 1.43 [1.29, 1.60] & $-$0.114 [$-$0.14, $-$0.09] & $-$0.218 & 0.037 & 12.2 \\
0.995 & 60 & char & def. & 4.5 & 3.6 & 0.5 & 3.8 & akbc & 1.19 [1.15, 1.23] & 1.24 [1.14, 1.35] & $-$0.092 [$-$0.12, $-$0.07] & $-$0.183 & 0.021 & 24.4 \\
0.995 & 60 & char & 5 & 4.8 & 4.3 & 0.5 & 3.7 & abkc & 1.29 [1.25, 1.34] & 1.10 [1.02, 1.19] & $-$0.083 [$-$0.11, $-$0.06] & $-$0.169 & 0.021 & 27.3 \\
0.995 & 60 & char & 30 & 4.4 & 3.5 & 0.5 & 3.8 & akbc & 1.15 [1.12, 1.19] & 1.26 [1.17, 1.37] & $-$0.091 [$-$0.12, $-$0.06] & $-$0.184 & 0.021 & 23.3 \\
10\,s & 15 & net & def. & 27.9 & 24.6 & 17.6 & 20.9 & abkc & 1.33 [1.30, 1.36] & 1.13 [1.06, 1.21] & $-$0.190 [$-$0.22, $-$0.16] & $-$0.196 & 0.085 & 0.4 \\
10\,s & 15 & net & 5 & 34.8 & 35.9 & 17.3 & 21.7 & bakc & 1.61 [1.58, 1.63] & 0.97 [0.94, 1.00] & $-$0.208 [$-$0.24, $-$0.18] & $-$0.223 & 0.090 & 0.7 \\
10\,s & 15 & net & 30 & 27.8 & 24.3 & 17.3 & 21.0 & abkc & 1.32 [1.30, 1.35] & 1.14 [1.07, 1.22] & $-$0.189 [$-$0.21, $-$0.16] & $-$0.192 & 0.085 & 0.3 \\
10\,s & 15 & char & def. & 22.4 & 20.7 & 14.5 & 15.5 & abkc & 1.45 [1.42, 1.48] & 1.08 [1.02, 1.14] & $-$0.196 [$-$0.22, $-$0.17] & $-$0.187 & 0.061 & 1.2 \\
10\,s & 15 & char & 5 & 27.7 & 29.8 & 14.0 & 16.0 & bakc & 1.73 [1.70, 1.76] & 0.93 [0.91, 0.95] & $-$0.210 [$-$0.24, $-$0.18] & $-$0.209 & 0.064 & 1.6 \\
10\,s & 15 & char & 30 & 22.3 & 20.4 & 14.0 & 15.6 & abkc & 1.43 [1.40, 1.46] & 1.09 [1.03, 1.16] & $-$0.194 [$-$0.22, $-$0.16] & $-$0.183 & 0.061 & 1.1 \\
10\,s & 30 & net & def. & 20.5 & 19.2 & 6.3 & 12.9 & abkc & 1.59 [1.55, 1.63] & 1.06 [1.01, 1.13] & $-$0.240 [$-$0.27, $-$0.21] & $-$0.265 & 0.048 & 2.6 \\
10\,s & 30 & net & 5 & 25.1 & 27.5 & 6.2 & 13.3 & bakc & 1.90 [1.85, 1.94] & 0.91 [0.89, 0.95] & $-$0.248 [$-$0.27, $-$0.22] & $-$0.277 & 0.049 & 3.5 \\
10\,s & 30 & net & 30 & 20.5 & 19.1 & 6.2 & 13.0 & abkc & 1.58 [1.54, 1.62] & 1.07 [1.01, 1.14] & $-$0.241 [$-$0.27, $-$0.21] & $-$0.264 & 0.046 & 2.3 \\
10\,s & 30 & char & def. & 14.9 & 14.8 & 4.4 & 8.6 & abkc & 1.73 [1.69, 1.77] & 1.00 [0.97, 1.05] & $-$0.218 [$-$0.24, $-$0.19] & $-$0.236 & 0.026 & 6.4 \\
10\,s & 30 & char & 5 & 17.8 & 20.6 & 4.2 & 8.8 & bakc & 2.03 [1.98, 2.08] & 0.86 [0.85, 0.88] & $-$0.221 [$-$0.25, $-$0.19] & $-$0.242 & 0.026 & 7.7 \\
10\,s & 30 & char & 30 & 14.8 & 14.6 & 4.2 & 8.7 & abkc & 1.71 [1.67, 1.75] & 1.01 [0.97, 1.06] & $-$0.217 [$-$0.24, $-$0.19] & $-$0.235 & 0.025 & 6.0 \\
10\,s & 60 & net & def. & 13.6 & 13.8 & 1.5 & 7.1 & bakc & 1.92 [1.85, 1.98] & 0.99 [0.93, 1.04] & $-$0.234 [$-$0.26, $-$0.21] & $-$0.278 & 0.044 & 11.5 \\
10\,s & 60 & net & 5 & 16.2 & 19.0 & 1.5 & 7.3 & bakc & 2.23 [2.16, 2.31] & 0.86 [0.82, 0.90] & $-$0.232 [$-$0.26, $-$0.21] & $-$0.270 & 0.041 & 13.7 \\
10\,s & 60 & net & 30 & 13.8 & 13.9 & 1.6 & 7.3 & bakc & 1.89 [1.83, 1.96] & 0.99 [0.94, 1.05] & $-$0.241 [$-$0.27, $-$0.21] & $-$0.281 & 0.040 & 10.7 \\
10\,s & 60 & char & def. & 8.5 & 9.2 & 0.8 & 4.1 & bakc & 2.07 [2.01, 2.13] & 0.92 [0.89, 0.95] & $-$0.187 [$-$0.21, $-$0.16] & $-$0.237 & 0.016 & 21.8 \\
10\,s & 60 & char & 5 & 9.5 & 11.8 & 0.7 & 4.1 & bakc & 2.34 [2.27, 2.42] & 0.80 [0.77, 0.84] & $-$0.176 [$-$0.20, $-$0.15] & $-$0.221 & 0.017 & 24.8 \\
10\,s & 60 & char & 30 & 8.5 & 9.2 & 0.7 & 4.2 & bakc & 2.03 [1.98, 2.09] & 0.93 [0.90, 0.96] & $-$0.188 [$-$0.22, $-$0.16] & $-$0.240 & 0.016 & 20.7 \\
30\,s & 15 & net & def. & 14.8 & 14.6 & 3.5 & 8.2 & abkc & 1.80 [1.76, 1.84] & 1.01 [0.97, 1.06] & $-$0.226 [$-$0.25, $-$0.20] & $-$0.244 & 0.025 & 7.1 \\
30\,s & 15 & net & 5 & 17.8 & 20.4 & 3.5 & 8.4 & bakc & 2.12 [2.07, 2.17] & 0.87 [0.85, 0.89] & $-$0.228 [$-$0.25, $-$0.20] & $-$0.248 & 0.025 & 8.5 \\
30\,s & 15 & net & 30 & 14.8 & 14.5 & 3.5 & 8.3 & abkc & 1.78 [1.73, 1.82] & 1.02 [0.98, 1.07] & $-$0.226 [$-$0.25, $-$0.20] & $-$0.243 & 0.024 & 6.7 \\
30\,s & 15 & char & def. & 13.3 & 13.5 & 3.2 & 7.4 & bakc & 1.79 [1.75, 1.84] & 0.99 [0.95, 1.02] & $-$0.216 [$-$0.24, $-$0.19] & $-$0.236 & 0.022 & 8.9 \\
30\,s & 15 & char & 5 & 15.8 & 18.5 & 3.1 & 7.5 & bakc & 2.09 [2.04, 2.15] & 0.85 [0.83, 0.87] & $-$0.216 [$-$0.24, $-$0.19] & $-$0.237 & 0.023 & 10.4 \\
30\,s & 15 & char & 30 & 13.3 & 13.3 & 3.1 & 7.5 & bakc & 1.77 [1.73, 1.81] & 0.99 [0.96, 1.03] & $-$0.216 [$-$0.24, $-$0.19] & $-$0.235 & 0.022 & 8.5 \\
30\,s & 30 & net & def. & 12.9 & 13.0 & 1.8 & 6.6 & bakc & 1.95 [1.90, 2.01] & 0.99 [0.94, 1.04] & $-$0.247 [$-$0.27, $-$0.22] & $-$0.280 & 0.022 & 10.6 \\
30\,s & 30 & net & 5 & 15.3 & 17.8 & 1.8 & 6.7 & bakc & 2.27 [2.21, 2.34] & 0.86 [0.83, 0.89] & $-$0.245 [$-$0.27, $-$0.22] & $-$0.276 & 0.023 & 12.6 \\
30\,s & 30 & net & 30 & 13.0 & 13.0 & 1.8 & 6.7 & bakc & 1.93 [1.87, 1.98] & 1.00 [0.95, 1.05] & $-$0.250 [$-$0.28, $-$0.22] & $-$0.282 & 0.020 & 9.9 \\
30\,s & 30 & char & def. & 10.0 & 10.6 & 1.3 & 5.0 & bakc & 1.97 [1.92, 2.03] & 0.94 [0.91, 0.97] & $-$0.202 [$-$0.23, $-$0.17] & $-$0.241 & 0.016 & 16.8 \\
30\,s & 30 & char & 5 & 11.4 & 13.9 & 1.3 & 5.0 & bakc & 2.27 [2.20, 2.33] & 0.82 [0.79, 0.85] & $-$0.195 [$-$0.22, $-$0.17] & $-$0.232 & 0.017 & 19.3 \\
30\,s & 30 & char & 30 & 10.0 & 10.5 & 1.3 & 5.1 & bakc & 1.95 [1.90, 2.00] & 0.95 [0.92, 0.99] & $-$0.203 [$-$0.23, $-$0.18] & $-$0.244 & 0.015 & 16.0 \\
30\,s & 60 & net & def. & 9.5 & 10.1 & 0.6 & 4.4 & bakc & 2.15 [2.09, 2.22] & 0.94 [0.89, 1.00] & $-$0.226 [$-$0.25, $-$0.20] & $-$0.285 & 0.018 & 20.9 \\
30\,s & 60 & net & 5 & 10.8 & 13.2 & 0.7 & 4.4 & bakc & 2.45 [2.36, 2.54] & 0.82 [0.78, 0.87] & $-$0.215 [$-$0.24, $-$0.19] & $-$0.264 & 0.017 & 24.3 \\
30\,s & 60 & net & 30 & 9.6 & 10.2 & 0.7 & 4.5 & bakc & 2.12 [2.05, 2.19] & 0.94 [0.89, 1.00] & $-$0.230 [$-$0.26, $-$0.20] & $-$0.286 & 0.017 & 19.7 \\
30\,s & 60 & char & def. & 6.5 & 7.4 & 0.3 & 3.0 & bakc & 2.17 [2.10, 2.25] & 0.88 [0.85, 0.92] & $-$0.160 [$-$0.19, $-$0.13] & $-$0.229 & 0.020 & 30.5 \\
30\,s & 60 & char & 5 & 6.9 & 8.9 & 0.3 & 2.9 & bakc & 2.40 [2.31, 2.49] & 0.77 [0.73, 0.82] & $-$0.143 [$-$0.17, $-$0.12] & $-$0.202 & 0.020 & 34.9 \\
30\,s & 60 & char & 30 & 6.5 & 7.3 & 0.3 & 3.1 & bakc & 2.13 [2.07, 2.20] & 0.89 [0.85, 0.93] & $-$0.160 [$-$0.19, $-$0.13] & $-$0.231 & 0.020 & 29.3 \\

\end{longtable}
}

{\scriptsize
\setlength{\tabcolsep}{3pt}
\begin{longtable}{@{}l l l r r r r l l r r@{}}
\caption{The pause family. Interval: between any two key presses or between two character presses; threshold in
seconds or as a writer quantile; Obs. as in Table~\ref{tab:s3}. Rates per hour; rank as in Table~\ref{tab:s3}. $\rho$:
score association with 95\% interval; adj.: given speed; len.: given essay and observation length.}\label{tab:spause}\\
\toprule
Interval & Threshold & Obs. & CS19 & CS21 & CoA & KL & Rank & $\rho$ & adj. & len. \\
\midrule
\endfirsthead
\toprule
Interval & Threshold & Obs. & CS19 & CS21 & CoA & KL & Rank & $\rho$ & adj. & len. \\
\midrule
\endhead
\bottomrule
\endfoot
% Generated from the aggregate result files by the authors' script; do not edit by hand.
any & 2\,s & def. & 143.8 & 105.5 & 193.3 & 179.8 & ckab & 0.025 [$-$0.00, 0.05] & 0.127 & 0.049 \\
any & 2\,s & 5 & 189.9 & 166.4 & 201.0 & 188.7 & cakb & $-$0.012 [$-$0.04, 0.02] & 0.083 & 0.056 \\
any & 2\,s & 30 & 146.6 & 108.0 & 200.8 & 180.2 & ckab & 0.019 [$-$0.01, 0.05] & 0.126 & 0.045 \\
any & 5\,s & def. & 69.7 & 56.0 & 93.0 & 76.9 & ckab & 0.000 [$-$0.03, 0.03] & 0.067 & 0.136 \\
any & 5\,s & 5 & 91.6 & 87.4 & 97.9 & 81.1 & cabk & $-$0.037 [$-$0.07, $-$0.01] & 0.020 & 0.147 \\
any & 5\,s & 30 & 71.4 & 57.7 & 97.8 & 77.6 & ckab & $-$0.011 [$-$0.04, 0.02] & 0.060 & 0.130 \\
any & 10\,s & def. & 42.5 & 36.6 & 50.8 & 38.8 & cakb & $-$0.086 [$-$0.11, $-$0.06] & $-$0.035 & 0.134 \\
any & 10\,s & 5 & 55.4 & 56.5 & 53.6 & 41.0 & back & $-$0.127 [$-$0.15, $-$0.10] & $-$0.087 & 0.139 \\
any & 10\,s & 30 & 43.7 & 37.9 & 53.5 & 39.4 & cakb & $-$0.101 [$-$0.13, $-$0.07] & $-$0.047 & 0.127 \\
any & 0.95 & def. & 95.0 & 51.2 & 489.1 & 377.9 & ckab & 0.616 [0.60, 0.63] & 0.413 & 0.074 \\
any & 0.95 & 5 & 123.7 & 80.0 & 498.3 & 393.7 & ckab & 0.583 [0.56, 0.60] & 0.347 & 0.081 \\
any & 0.95 & 30 & 94.8 & 51.1 & 497.8 & 375.2 & ckab & 0.619 [0.60, 0.64] & 0.420 & 0.073 \\
any & 0.99 & def. & 18.8 & 10.1 & 95.8 & 74.8 & ckab & 0.617 [0.60, 0.64] & 0.415 & 0.075 \\
any & 0.99 & 5 & 24.5 & 15.9 & 97.6 & 77.9 & ckab & 0.586 [0.57, 0.61] & 0.350 & 0.082 \\
any & 0.99 & 30 & 18.8 & 10.1 & 97.5 & 74.3 & ckab & 0.620 [0.60, 0.64] & 0.422 & 0.074 \\
any & 0.995 & def. & 9.3 & 5.0 & 46.6 & 36.8 & ckab & 0.619 [0.60, 0.64] & 0.418 & 0.075 \\
any & 0.995 & 5 & 12.1 & 7.8 & 47.4 & 38.4 & ckab & 0.588 [0.57, 0.61] & 0.354 & 0.081 \\
any & 0.995 & 30 & 9.2 & 5.0 & 47.4 & 36.6 & ckab & 0.621 [0.60, 0.64] & 0.423 & 0.073 \\
char & 2\,s & def. & 138.7 & 92.7 & 206.1 & 203.9 & ckab & 0.049 [0.02, 0.08] & 0.136 & 0.006 \\
char & 2\,s & 5 & 182.4 & 145.1 & 213.7 & 213.6 & ckab & 0.011 [$-$0.02, 0.04] & 0.087 & 0.008 \\
char & 2\,s & 30 & 140.8 & 94.2 & 213.5 & 203.9 & ckab & 0.045 [0.02, 0.07] & 0.137 & 0.002 \\
char & 5\,s & def. & 67.9 & 50.1 & 100.2 & 85.8 & ckab & 0.032 [0.00, 0.06] & 0.084 & 0.120 \\
char & 5\,s & 5 & 89.1 & 77.9 & 105.2 & 90.2 & ckab & $-$0.003 [$-$0.03, 0.02] & 0.041 & 0.130 \\
char & 5\,s & 30 & 69.4 & 51.4 & 105.1 & 86.3 & ckab & 0.021 [$-$0.01, 0.05] & 0.079 & 0.115 \\
char & 10\,s & def. & 41.8 & 33.6 & 55.2 & 42.2 & ckab & $-$0.058 [$-$0.09, $-$0.03] & $-$0.011 & 0.134 \\
char & 10\,s & 5 & 54.6 & 51.8 & 58.2 & 44.6 & cabk & $-$0.099 [$-$0.13, $-$0.07] & $-$0.061 & 0.139 \\
char & 10\,s & 30 & 43.1 & 34.8 & 58.2 & 42.8 & cakb & $-$0.074 [$-$0.10, $-$0.05] & $-$0.022 & 0.126 \\
char & 0.95 & def. & 75.0 & 35.5 & 443.5 & 305.5 & ckab & 0.660 [0.64, 0.68] & 0.500 & 0.058 \\
char & 0.95 & 5 & 97.6 & 55.1 & 451.4 & 318.0 & ckab & 0.620 [0.60, 0.64] & 0.429 & 0.061 \\
char & 0.95 & 30 & 74.8 & 35.2 & 450.9 & 303.0 & ckab & 0.665 [0.65, 0.68] & 0.508 & 0.057 \\
char & 0.99 & def. & 14.8 & 7.0 & 86.7 & 60.2 & ckab & 0.662 [0.65, 0.68] & 0.503 & 0.060 \\
char & 0.99 & 5 & 19.3 & 10.9 & 88.2 & 62.7 & ckab & 0.623 [0.61, 0.64] & 0.433 & 0.062 \\
char & 0.99 & 30 & 14.8 & 7.0 & 88.1 & 59.8 & ckab & 0.667 [0.65, 0.68] & 0.511 & 0.058 \\
char & 0.995 & def. & 7.3 & 3.4 & 42.0 & 29.5 & ckab & 0.663 [0.65, 0.68] & 0.504 & 0.057 \\
char & 0.995 & 5 & 9.5 & 5.3 & 42.7 & 30.8 & ckab & 0.626 [0.61, 0.64] & 0.438 & 0.060 \\
char & 0.995 & 30 & 7.2 & 3.4 & 42.7 & 29.3 & ckab & 0.667 [0.65, 0.68] & 0.512 & 0.056 \\

\end{longtable}
}

\section{Overdispersion}
\label{app:phi}

The threshold of 1.5 was set in the primary conditions as the value at which we would replace a Poisson count model with
a negative binomial one; it is a working threshold, not a published rule, which is why the main text reports the
distribution of $\varphi$ and the count of definitions whose interval exceeds 1. With a single pooled rate as the
expectation, the baseline $\varphi$ was larger at every level and also defined at the level of writers
(Table~\ref{tab:s4b}); with the writer's own rate, CoAuthor sessions could not be told apart from counting.
Table~\ref{tab:s4} gives $\varphi$ with the writer's rate for every definition. Writers with no stall carry no
information about their own rate and are left out of the writer-expectation $\varphi$; their share differs between
definitions (in KLiCKe 0.1--35\%), so the definitions differ in the writers they use.

\begin{center}\footnotesize
\begin{tabular}{@{}l l r r@{}}
\toprule
Dataset & Level & Pooled rate & Writer's rate \\
\midrule
CS1 2019 & block & 1.97 [1.85, 2.10] & 1.44 [1.35, 1.53] \\
 & day & 2.25 [2.08, 2.41] & 1.57 [1.45, 1.70] \\
 & assignment & 2.90 [2.64, 3.17] & 1.95 [1.76, 2.15] \\
 & writer & 4.91 [4.23, 5.57] & \\
CS1 2021 & block & 2.26 [1.88, 2.66] & 1.82 [1.57, 2.10] \\
 & day & 3.13 [2.54, 3.74] & 2.44 [2.05, 2.90] \\
 & assignment & 5.12 [3.96, 6.30] & 3.82 [3.11, 4.59] \\
 & writer & 10.18 [5.58, 15.15] & \\
CoAuthor & session & 1.76 [1.29, 2.20] & 1.08 [0.98, 1.19] \\
 & writer & 20.98 [8.44, 35.11] & \\
KLiCKe & writer & 2.57 [2.34, 2.81] & \\
\bottomrule
\end{tabular}
\captionof{table}{Baseline $\varphi$ with two expectations, with 95\% intervals.}\label{tab:s4b}
\end{center}

{\scriptsize
\begin{longtable}{@{}l r l l r r r r r@{}}
\caption{$\varphi$ with the writer's rate as expectation, for every definition. Columns: CS1 2019 blocks and
assignments, CS1 2021 blocks and assignments, CoAuthor sessions.}\label{tab:s4}\\
\toprule
$\theta$ & $W$ & A & Obs. & CS19 bl. & CS19 as. & CS21 bl. & CS21 as. & CoA ses. \\
\midrule
\endfirsthead
\toprule
$\theta$ & $W$ & A & Obs. & CS19 bl. & CS19 as. & CS21 bl. & CS21 as. & CoA ses. \\
\midrule
\endhead
\bottomrule
\endfoot
% Generated from the aggregate result files by the authors' script; do not edit by hand.
0.95 & 15 & net & def. & 3.26 & 3.64 & 4.64 & 7.71 & 0.93 \\
0.95 & 15 & net & 5 & 1.37 & 1.54 & 1.59 & 2.13 & 0.93 \\
0.95 & 15 & net & 30 & 3.28 & 3.72 & 4.66 & 7.87 & 0.93 \\
0.95 & 15 & char & def. & 2.19 & 2.54 & 3.36 & 5.43 & 0.80 \\
0.95 & 15 & char & 5 & 0.91 & 1.05 & 1.04 & 1.06 & 0.80 \\
0.95 & 15 & char & 30 & 2.21 & 2.59 & 3.33 & 5.37 & 0.80 \\
0.95 & 30 & net & def. & 2.42 & 2.55 & 3.26 & 5.18 & 1.19 \\
0.95 & 30 & net & 5 & 1.32 & 1.38 & 1.18 & 1.39 & 1.17 \\
0.95 & 30 & net & 30 & 2.46 & 2.62 & 3.30 & 5.55 & 1.17 \\
0.95 & 30 & char & def. & 1.30 & 1.41 & 2.05 & 3.21 & 0.97 \\
0.95 & 30 & char & 5 & 0.72 & 0.79 & 0.74 & 0.71 & 0.95 \\
0.95 & 30 & char & 30 & 1.31 & 1.42 & 2.02 & 3.18 & 0.95 \\
0.95 & 60 & net & def. & 1.92 & 1.99 & 2.16 & 3.25 & 1.36 \\
0.95 & 60 & net & 5 & 1.45 & 1.53 & 1.07 & 1.22 & 1.37 \\
0.95 & 60 & net & 30 & 1.95 & 2.05 & 2.24 & 3.69 & 1.37 \\
0.95 & 60 & char & def. & 0.87 & 0.90 & 1.09 & 1.47 & 0.98 \\
0.95 & 60 & char & 5 & 0.77 & 0.79 & 0.67 & 0.74 & 0.97 \\
0.95 & 60 & char & 30 & 0.87 & 0.90 & 1.08 & 1.50 & 0.97 \\
0.99 & 15 & net & def. & 1.97 & 3.03 & 2.35 & 5.18 & 1.00 \\
0.99 & 15 & net & 5 & 1.17 & 1.77 & 1.28 & 2.82 & 1.00 \\
0.99 & 15 & net & 30 & 1.97 & 3.07 & 2.31 & 5.33 & 1.00 \\
0.99 & 15 & char & def. & 1.48 & 2.11 & 1.96 & 4.26 & 0.94 \\
0.99 & 15 & char & 5 & 0.87 & 1.19 & 1.03 & 1.96 & 0.95 \\
0.99 & 15 & char & 30 & 1.46 & 2.16 & 1.89 & 4.32 & 0.95 \\
0.99$^\ast$ & 30 & net & def. & 1.44 & 1.95 & 1.82 & 3.82 & 1.08 \\
0.99 & 30 & net & 5 & 0.95 & 1.21 & 1.08 & 2.08 & 1.04 \\
0.99 & 30 & net & 30 & 1.46 & 2.00 & 1.84 & 4.09 & 1.04 \\
0.99 & 30 & char & def. & 1.02 & 1.27 & 1.31 & 2.54 & 0.99 \\
0.99 & 30 & char & 5 & 0.75 & 0.87 & 0.82 & 1.28 & 0.99 \\
0.99 & 30 & char & 30 & 1.02 & 1.28 & 1.30 & 2.63 & 0.99 \\
0.99 & 60 & net & def. & 1.09 & 1.28 & 1.33 & 2.44 & 1.18 \\
0.99 & 60 & net & 5 & 0.92 & 1.02 & 0.93 & 1.36 & 1.09 \\
0.99 & 60 & net & 30 & 1.09 & 1.31 & 1.38 & 2.71 & 1.09 \\
0.99 & 60 & char & def. & 0.78 & 0.83 & 0.90 & 1.28 & 1.00 \\
0.99 & 60 & char & 5 & 0.78 & 0.77 & 0.78 & 0.99 & 1.00 \\
0.99 & 60 & char & 30 & 0.77 & 0.82 & 0.89 & 1.34 & 1.00 \\
0.995 & 15 & net & def. & 1.64 & 2.67 & 1.64 & 3.94 & 1.05 \\
0.995 & 15 & net & 5 & 1.17 & 1.71 & 1.26 & 2.69 & 1.03 \\
0.995 & 15 & net & 30 & 1.67 & 2.73 & 1.69 & 4.23 & 1.03 \\
0.995 & 15 & char & def. & 1.32 & 2.01 & 1.46 & 3.35 & 1.01 \\
0.995 & 15 & char & 5 & 0.97 & 1.30 & 1.11 & 2.22 & 1.02 \\
0.995 & 15 & char & 30 & 1.33 & 2.04 & 1.48 & 3.55 & 1.02 \\
0.995 & 30 & net & def. & 1.27 & 1.83 & 1.42 & 3.24 & 1.10 \\
0.995 & 30 & net & 5 & 1.01 & 1.27 & 1.13 & 2.19 & 1.08 \\
0.995 & 30 & net & 30 & 1.28 & 1.86 & 1.46 & 3.52 & 1.08 \\
0.995 & 30 & char & def. & 0.96 & 1.25 & 1.13 & 2.35 & 1.07 \\
0.995 & 30 & char & 5 & 0.85 & 0.94 & 0.97 & 1.53 & 1.06 \\
0.995 & 30 & char & 30 & 0.97 & 1.27 & 1.14 & 2.49 & 1.06 \\
0.995 & 60 & net & def. & 1.00 & 1.26 & 1.11 & 2.31 & 1.06 \\
0.995 & 60 & net & 5 & 0.94 & 1.03 & 1.08 & 1.85 & 1.04 \\
0.995 & 60 & net & 30 & 1.02 & 1.30 & 1.21 & 2.62 & 1.04 \\
0.995 & 60 & char & def. & 0.76 & 0.83 & 0.85 & 1.37 & 1.01 \\
0.995 & 60 & char & 5 & 0.81 & 0.80 & 0.96 & 1.24 & 0.99 \\
0.995 & 60 & char & 30 & 0.76 & 0.85 & 0.87 & 1.51 & 0.99 \\
10\,s & 15 & net & def. & 2.27 & 2.44 & 3.36 & 5.76 & 0.94 \\
10\,s & 15 & net & 5 & 1.04 & 1.18 & 0.91 & 0.97 & 0.96 \\
10\,s & 15 & net & 30 & 2.31 & 2.49 & 3.39 & 5.90 & 0.96 \\
10\,s & 15 & char & def. & 1.68 & 1.85 & 2.61 & 4.33 & 0.99 \\
10\,s & 15 & char & 5 & 0.77 & 0.89 & 0.69 & 0.67 & 1.01 \\
10\,s & 15 & char & 30 & 1.71 & 1.88 & 2.61 & 4.32 & 1.01 \\
10\,s & 30 & net & def. & 1.78 & 1.97 & 2.42 & 4.00 & 1.06 \\
10\,s & 30 & net & 5 & 1.02 & 1.20 & 0.81 & 0.91 & 1.08 \\
10\,s & 30 & net & 30 & 1.82 & 2.01 & 2.51 & 4.23 & 1.08 \\
10\,s & 30 & char & def. & 1.09 & 1.20 & 1.58 & 2.55 & 1.02 \\
10\,s & 30 & char & 5 & 0.67 & 0.78 & 0.56 & 0.63 & 1.02 \\
10\,s & 30 & char & 30 & 1.10 & 1.20 & 1.61 & 2.55 & 1.02 \\
10\,s & 60 & net & def. & 1.49 & 1.63 & 1.56 & 2.49 & 1.12 \\
10\,s & 60 & net & 5 & 1.20 & 1.38 & 0.79 & 0.89 & 1.11 \\
10\,s & 60 & net & 30 & 1.51 & 1.67 & 1.66 & 2.77 & 1.11 \\
10\,s & 60 & char & def. & 0.79 & 0.82 & 0.87 & 1.27 & 0.99 \\
10\,s & 60 & char & 5 & 0.75 & 0.79 & 0.58 & 0.76 & 1.00 \\
10\,s & 60 & char & 30 & 0.79 & 0.81 & 0.86 & 1.26 & 1.00 \\
30\,s & 15 & net & def. & 1.15 & 1.23 & 1.57 & 2.56 & 1.02 \\
30\,s & 15 & net & 5 & 0.75 & 0.84 & 0.57 & 0.66 & 1.02 \\
30\,s & 15 & net & 30 & 1.17 & 1.25 & 1.61 & 2.62 & 1.02 \\
30\,s & 15 & char & def. & 0.99 & 1.07 & 1.39 & 2.28 & 1.06 \\
30\,s & 15 & char & 5 & 0.68 & 0.77 & 0.56 & 0.66 & 1.04 \\
30\,s & 15 & char & 30 & 1.00 & 1.08 & 1.42 & 2.26 & 1.04 \\
30\,s & 30 & net & def. & 1.15 & 1.19 & 1.34 & 2.15 & 1.06 \\
30\,s & 30 & net & 5 & 0.88 & 0.97 & 0.61 & 0.77 & 1.06 \\
30\,s & 30 & net & 30 & 1.17 & 1.22 & 1.40 & 2.33 & 1.06 \\
30\,s & 30 & char & def. & 0.84 & 0.88 & 1.01 & 1.55 & 1.00 \\
30\,s & 30 & char & 5 & 0.70 & 0.77 & 0.58 & 0.76 & 1.00 \\
30\,s & 30 & char & 30 & 0.84 & 0.88 & 1.02 & 1.57 & 1.00 \\
30\,s & 60 & net & def. & 1.01 & 1.09 & 1.04 & 1.48 & 1.08 \\
30\,s & 60 & net & 5 & 0.97 & 1.10 & 0.66 & 0.70 & 1.08 \\
30\,s & 60 & net & 30 & 1.03 & 1.11 & 1.09 & 1.59 & 1.08 \\
30\,s & 60 & char & def. & 0.72 & 0.74 & 0.73 & 0.93 & 1.04 \\
30\,s & 60 & char & 5 & 0.78 & 0.82 & 0.63 & 0.78 & 1.06 \\
30\,s & 60 & char & 30 & 0.72 & 0.74 & 0.73 & 0.94 & 1.06 \\

\end{longtable}
}

\section{The length model}
\label{app:length}

The model is $\log L_{i} = \mu + w_{a} + u_{b} + d_{c} + k_{e} + \varepsilon_{i}$ for stall $i$ of writer $a$, assignment
$b$, day $c$ and work block $e$ in CS1, and $\log L_{i} = \mu + w_{a} + v_{b} + \varepsilon_{i}$ with session $b$ in
CoAuthor, all effects normal and independent, fitted by REML with the writer as group and the lower levels as variance
components. The default optimiser did not converge in 14 of the 27 fits of the primary computation and, for the
baseline CS1 2019, reached a likelihood 241 units lower than Powell and Nelder--Mead, which agreed with each other
(shares within 0.01). We therefore took the converged solution with the largest REML likelihood; all 270 fits over the
90 definitions converged. On data simulated with the unit structure and known components, the fit recovered all levels
within $\pm 30\%$ in CS1 2019 and CoAuthor; in CS1 2021 (31 writers) the assignment and day components were
underestimated in one of two simulated sets. Stalls that did not resume within their observation interval (under the
baseline 1,232 in CS1 2019, 383 in CS1 2021 and 57 in CoAuthor) were excluded. In CoAuthor the three definitions that
failed (window 60~s, quantile 0.995, characters only, with each segmentation) had 86--88 stalls and session shares of
0.52--0.58 against stall shares of 0.42--0.48.

\section{The count decomposition}
\label{app:counts}

A Poisson log-normal mixed model with writer, assignment, day and block (CS1) or writer and session (CoAuthor),
estimated by variational Bayes, did not recover simulated components: in CS1 2019 no level was recovered within
$\pm 30\%$ in either simulated set, and in CS1 2021 only the writer and the shared assignment effect were. On the real
data the shares were unstable across definitions, and repeated runs on the same input with the same seed moved a share
by up to 0.93. With on average 2.02 blocks per assignment in CS1 2019 and 3.55 in CS1 2021, this estimator cannot
separate the levels of count variation. We report no shares.

\section{Scope of the search for prior work}
\label{app:search}

Four rounds of searches (OpenAlex, Crossref, publisher pages and full texts where available, between 2026-11-05 and
2026-11-08) covered analytic choices and multiverse analyses in learning analytics, educational measurement and
behavioural data; time-on-task and pause thresholds in CS1 and writing, with the answers to them (mixture models,
personalised thresholds, state models); decompositions and dispersion of keystroke measures; and associations of pause
measures with essay scores, with adjustment for speed and length and with null, positive or threshold-dependent
results. Statements that we found no study of a given
kind are limited to these searches.

\section{Relation to the companion study}
\label{app:companion}

\begin{center}\footnotesize
\begin{tabularx}{\textwidth}{@{}l X X@{}}
\toprule
 & Companion study \cite{miura2026runup} & This study \\
\midrule
Question & Do the 30-s changes of an estimator's signals distinguish the run-up to a stall from a matched short
pause? & Which conclusions drawn from a stall label move across definitions? \\
Design & Preregistered equivalence test on a holdout & Multiverse of 89 definitions with writer-level intervals \\
Main result & Both signals inside the equivalence band & Score association, rate ratios and overdispersion move;
CS1 length structure stays \\
Data & CS1 2019 (holdout), KLiCKe and CoAuthor as descriptions & CS1 2019 and 2021, KLiCKe, CoAuthor \\
Shared & \multicolumn{2}{X}{The baseline stall definition, the conversion of CS1, the estimator, the opened holdouts
and the data.} \\
New here & \multicolumn{2}{X}{The 89 definitions, the pause family, the intervals, the score association and its
adjustments, the overdispersion within writers and the length decomposition.} \\
\bottomrule
\end{tabularx}
\end{center}

Some baseline values (rates and part of $\varphi$) can coincide with values reported in the companion study, because
the definition is the same. Here the baseline is used only as a reference point.
